\section{Methodology}
\label{sec:methodology}

\subsection{Data and Target}

The review corpus and film plots serve separate purposes: reviews define the control target, while plots provide generation inputs. Because reviews have no movie identifiers, generated responses are not treated as predictions of film-level audience behaviour.

The 5,000-row review corpus was audited to a frozen 4,625-row split surface~\cite{b4}: Gold-300 contains 300 items, R1 contains 2,162, and R2 contains 2,163. Gold-300 never enters training, retrieval, prompting, or threshold selection. LaBSE-based diagnostics~\cite{b3} did not support discrete audience personas. Two native-Bangla annotators instead tested a two-level engagement-specificity continuum through length-matched intrusion and direction judgments; the numerical evidence is reported with the results.

\subsection{Verifier-Isolated Framework}

The framework uses four functional roles in a fixed loop, shown in Figure~\ref{fig:methodology-workflow}. Researcher and Critic are deterministic tool roles; only Writer and Reflector call the generator.

\begin{figure*}[t]
\centering
\includegraphics[width=\textwidth]{methodology_workflow.png}
\caption{Bounded generation and evaluator isolation. Solid arrows show the generation path; dashed arrows show revision and outcome-only evaluation. Verifier-B cannot affect retrieval, revision, or stopping.}
\label{fig:methodology-workflow}
\end{figure*}

\begin{table}[t]
\centering
\caption{Data roles and registered workflow settings}
\label{tab:workflow-settings}
\footnotesize
\setlength{\tabcolsep}{3.5pt}
\renewcommand{\arraystretch}{1.04}
\begin{tabularx}{\columnwidth}{@{}
  >{\raggedright\arraybackslash}p{0.25\columnwidth}
  >{\raggedright\arraybackslash}X@{}}
\toprule
\textbf{Component} & \textbf{Registered setting} \\
\midrule
Target & Level~0: general or formulaic; Level~1: engagement with a specific narrative detail \\
Construct data & Gold-300: human evaluation only; R1: Verifier-A and retrieval; R2: Verifier-B only \\
Retrieval & 886 Region-A R1 reviews; ten same-level LaBSE neighbours~\cite{b8} \\
Acceptance & Frozen LaBSE logistic Verifier-A; maximum three Writer attempts \\
Diagnosis & Deterministic length, orthographic, connective, sentiment-marker, and lexical-richness rules \\
Outcome scoring & BanglaBERT Verifier-B~\cite{b2}, trained on 888 disjoint Region-A R2 reviews and applied only to sealed outputs \\
Trace & Retrieved IDs, drafts, scores, feedback, tokens, attempts, and stopping reason \\
\bottomrule
\end{tabularx}
\end{table}

Verifier-A alone controls acceptance in the main neuro-symbolic condition; symbolic rules only localize revision feedback. Retrieved reviews provide style and specificity examples, not facts about the input plot. A common 20-word ceiling limits output length, but length diagnostics remain necessary because the generated levels are not length-neutral.

\subsection{Experimental Design}

The frozen experiment crosses 90 held-out plots, two requested levels, ten conditions, and three paired seeds (42--44), producing 5,400 outputs. The conditions are zero-shot, static few-shot, RAG-only, neural-gated RAG, symbolic-gated RAG, neural-gated RAG with symbolic feedback, intrinsic self-critique, external-role self-critique, a hosted Gemma-4 judge loop, and compute-matched blind resampling. Every condition contributes 540 cases.

Verifier-B target probability is the primary outcome. Each active condition is paired with zero-shot by plot, level, and seed. We use 10,000 paired-bootstrap resamples, 95\% confidence intervals, McNemar tests for binary target match, and Benjamini--Hochberg correction across the nine registered comparisons~\cite{b49,b50}. The seeds are pairing blocks, not independent studies. A separate balanced subset of 100 outputs was judged by three blinded native-Bangla annotators. The planned tests compare active conditions with zero-shot; they do not rank the active systems.
